\begin{ouverture}[calcul]{1}{Les nombres jusqu’au million}{Partie 1 · Nombres entiers et calculs}
{Classes · Rangs · Décompositions · Demi-droite}
\objectif{Lire, écrire et décomposer les nombres jusqu’à six chiffres.}
\objectif{Distinguer le chiffre des dizaines et le nombre de dizaines.}
\objectif{Comparer, encadrer et placer des nombres entiers.}
\prerequis{Numération jusqu’à 9 999 et tables d’addition.}
\end{ouverture}
\pagestyle{otMain}\essentiel\fontsize{11}{15}\selectfont
\phantomsection\label{cours-1-1}
\essSection{1}{Lire et décomposer les grands nombres}
\begin{definition}Dans un nombre entier, les chiffres sont groupés par classes de trois rangs : la classe des unités simples, puis la classe des mille. Le nombre $304\,065$ se lit « trois-cent-quatre-mille-soixante-cinq » : 304 dans la classe des mille, 065 dans la classe des unités.\end{definition}\begin{center}\renewcommand{\arraystretch}{1.35}\begin{tabular}{c|cc}\toprule
Classe & des mille & des unités simples\\\midrule
Rangs & centaines, dizaines, unités & centaines, dizaines, unités\\
$304\,065$ & 3 0 4 & 0 6 5\\
\bottomrule\end{tabular}\end{center}
\begin{propriete}La valeur d’un chiffre dépend de sa position : dans $52\,700$, le chiffre 2 vaut $2\,000$. Dix unités font une dizaine, dix dizaines font une centaine, dix centaines font un millier : chaque rang vaut dix fois le rang placé à sa droite.\end{propriete}
\begin{exemple}$304\,065=3\times100\,000+4\times1\,000+6\times10+5$. Les zéros gardent la place des rangs vides : sans eux, on lirait $3\,465$, un tout autre nombre.\end{exemple}
\begin{savoirfaire}[Chiffre ou nombre]Dans $4\,627$, le chiffre des dizaines est 2. Le nombre de dizaines est 462, car $4\,627=462\times10+7$ : on compte toutes les dizaines contenues dans le nombre. Cette distinction sert pour les échanges : 462 carnets de dix timbres sont complets.\end{savoirfaire}
\clearpage
\phantomsection\label{cours-1-2}
\essSection{2}{Comparer, encadrer et placer}
\begin{propriete}Pour comparer deux entiers : celui qui a le moins de chiffres est le plus petit ; à nombre de chiffres égal, on compare rang par rang en partant de la gauche. On écrit $9\,875<12\,034$ et $64\,350>64\,280$.\end{propriete}
\label{fig-cours-1-2-cours-droite-mille}\figureGeo{cours-droite-mille}

\begin{savoirfaire}[Encadrer et intercaler]Encadrer $45\,670$ entre deux dizaines de mille consécutives : $40\,000<45\,670<50\,000$. Intercaler un nombre entre $23\,500$ et $23\,600$ : par exemple $23\,540$. « Compris entre » se vérifie avec les deux symboles.\end{savoirfaire}
\begin{attention}Sur une demi-droite graduée, lire d’abord la valeur d’un intervalle entre deux graduations. Un nombre n’est pas plus grand parce que son écriture semble plus longue à l’œil : comparer les rangs, pas la place occupée sur la page.\end{attention}

\clearpage\rubrique{Découverte}
\begin{blocLivre}{ex-1-1}\begin{exercice}[1]{Écrire en chiffres}
Écrire en chiffres : « huit-mille-quarante » ; « vingt-trois-mille-quatre-cent-cinq » ; « six-cent-mille-neuf ».
\end{exercice}\end{blocLivre}
\begin{blocLivre}{ex-1-2}\begin{exercice}[1]{Lire en classes}
Écrire en lettres les nombres $86\,312$ et $450\,070$.
\end{exercice}\end{blocLivre}
\begin{blocLivre}{ex-1-3}\begin{exercice}[1]{La valeur d’un chiffre}
Dans $582\,741$, donner le chiffre des unités de mille, puis la valeur du chiffre 8.
\end{exercice}\end{blocLivre}
\begin{blocLivre}{ex-1-4}\begin{exercice}[1]{Chiffre et nombre de dizaines}
Pour le nombre $4\,627$, donner le chiffre des dizaines puis le nombre de dizaines.
\end{exercice}\end{blocLivre}
\begin{blocLivre}{ex-1-5}\begin{exercice}[1]{Comparer deux nombres}
Compléter par $<$ ou $>$ : $45\,806\ \ldots\ 45\,860$, puis $9\,875\ \ldots\ 12\,034$.
\end{exercice}\end{blocLivre}

\rubrique{Entraînement}
\begin{blocLivre}{ex-1-6}\begin{exercice}[2]{Décomposer}
Écrire la décomposition de $304\,065$ avec les valeurs de chaque chiffre.
\end{exercice}\end{blocLivre}
\begin{blocLivre}{ex-1-7}\begin{exercice}[2]{Recomposer}
Quel nombre s’écrit $7\times10\,000+3\times100+9$ ? Et $5\times100\,000+2\times10$ ?
\end{exercice}\end{blocLivre}
\begin{blocLivre}{ex-1-8}\begin{exercice}[2]{Sur la demi-droite}
Lire les nombres repérés par les points A, B et C.
\label{fig-ex-1-8}\figureGeo{num-8}
\end{exercice}\end{blocLivre}
\begin{blocLivre}{ex-1-9}\begin{exercice}[2]{Encadrer}
Encadrer $45\,670$ entre deux dizaines de mille consécutives, puis $8\,294$ entre deux centaines consécutives.
\end{exercice}\end{blocLivre}
\begin{blocLivre}{ex-1-10}\begin{exercice}[2]{Ordonner}
Ranger dans l’ordre croissant : $56\,400$ ; $9\,980$ ; $56\,040$ ; $100\,000$ ; $55\,999$.
\end{exercice}\end{blocLivre}

\rubrique{Problèmes}
\begin{blocLivre}{ex-1-11}\begin{exercice}[3]{Des fleuves à ranger}
La Loire mesure $1\,006$ km, la Seine 775 km, la Garonne 647 km et le Rhône 812 km. Ranger ces fleuves du plus court au plus long. Lesquels mesurent moins de 800 km ?
\end{exercice}\end{blocLivre}
\begin{blocLivre}{ex-1-12}\begin{exercice}[3]{Des carnets de timbres}
Une collection compte $3\,240$ timbres, rangés par carnets de 10. Combien de carnets complets obtient-on ? Même question pour $5\,278$ timbres, puis donner le nombre de timbres restants.
\end{exercice}\end{blocLivre}
\begin{blocLivre}{ex-1-13}\begin{exercice}[3]{Un nombre à inventer}
Trouver un nombre compris entre $45\,300$ et $45\,400$ dont le chiffre des dizaines est 7. Combien y a-t-il de possibilités ?
\end{exercice}\end{blocLivre}
\begin{blocLivre}{ex-1-14}\begin{exercice}[3]{Avec cinq cartes}
Avec les cartes 7, 0, 3, 9 et 1 utilisées chacune une fois, former le plus grand nombre possible, puis le plus petit nombre à cinq chiffres.
\end{exercice}\end{blocLivre}
\begin{blocLivre}{ex-1-15}\begin{exercice}[3]{Trois communes}
Trois communes comptent $45\,806$, $45\,860$ et $9\,980$ habitants. Ranger les populations dans l’ordre décroissant, encadrer la plus grande entre deux dizaines de mille consécutives et donner son chiffre des centaines.
\end{exercice}\end{blocLivre}

\clearpage
